% Generated draft: requires mathematical and source review.
\documentclass[11pt]{article}
\usepackage[margin=25mm]{geometry}
\usepackage{fontspec}
\setmainfont{Latin Modern Roman}
\usepackage{amsmath,amssymb,mathtools,unicode-math}
\setmathfont{Latin Modern Math}
\usepackage{enumitem,needspace,xurl,hyperref}
\setlength{\parindent}{0pt}
\setlength{\parskip}{5pt}
\setlength{\emergencystretch}{3em}
\newcommand{\N}{\mathbb{N}}
\newcommand{\Z}{\mathbb{Z}}
\newcommand{\Q}{\mathbb{Q}}
\newcommand{\R}{\mathbb{R}}
\newcommand{\C}{\mathbb{C}}
\newcommand{\F}{\mathbb{F}}
\newcommand{\A}{\mathbb{A}}
\newcommand{\PP}{\mathbb{P}}
\newcommand{\E}{\mathbb{E}}
\newcommand{\eps}{\varepsilon}
\newcommand{\dd}{\,\mathrm{d}}
\begin{document}
\section*{UnsolvedMath ID: 700999}
\textbf{Source problem code:} catalogue entry
\subsection{Definitions and mathematical statement}
Let \(F_2=\langle a,b\rangle\) be the free group of rank two. For \(x,y\in F_2\), put \([x,y]=x^{-1}y^{-1}xy\), and let \(F_2'=[F_2,F_2]\) be its derived subgroup. For \(g\in F_2'\), define its commutator length by
\[
\operatorname{cl}_{F_2}(g)=\min\left\{r\in\Z_{\geq 0}:g=\prod_{i=1}^{r}[x_i,y_i],\quad x_i,y_i\in F_2\right\}.
\]
The empty product is the identity, so \(\operatorname{cl}_{F_2}(1)=0\).

Let \(\mathcal N\) denote the natural numbers in the source's convention. The supplied record does not specify whether \(0\in\mathcal N\); this boundary convention remains unresolved. A pair \((k,m)\in\mathcal N^2\) is called \emph{admissible} if
\[
\exists w\in F_2'\qquad \operatorname{cl}_{F_2}(w^m)=k.
\]
The problem is to determine exactly the set
\[
\left\{(k,m)\in\mathcal N^2:\exists w\in F_2',\ \operatorname{cl}_{F_2}(w^m)=k\right\}.
\]
A complete formulation must specify whether \(\mathcal N=\Z_{\geq1}\) or \(\mathcal N=\Z_{\geq0}\).

The supplied source reports the following known partial result, attributed to V. G. Bardakov: for every integer \(k\geq2\), the pair \((k,2)\) is admissible. This is retained as a source-reported result, without independent verification.
\subsection{Short English statement}
Determine which pairs of natural numbers \((k,m)\) occur as the commutator length \(k\) of an \(m\)-th power of an element of the derived subgroup of the free group of rank two. The source leaves the inclusion of zero among natural numbers unspecified and reports that every \((k,2)\) with \(k\geq2\) is admissible.
\subsection{Sources}
\begin{itemize}
\item Local UnsolvedMath registry, permanent ID \(700999\), exact source code \texttt{catalogue entry}, title \emph{Kourovka Notebook Problem 14.14}; supplied upstream revision \texttt{d371ee917ac9647fbb06cd6fc529deec6f3dda82}. The registry labels the problem \emph{open}; this is a supplied classification, not an independently verified assessment.
\item \emph{The Kourovka Notebook (Unsolved Problems in Group Theory)}, original problem 14.14, 14th issue (1999); the supplied background identifies the consulted edition as the 21st issue edition (2026), \texttt{arXiv:1401.0300v48}. The collection description names E. I. Khukhro and V. D. Mazurov as editors. These bibliographic details are supplied and have not been independently checked.
\item Attribution: the statement begins with the names C. C. Edmunds and G. Rosenberger; the background and proposer field instead identify V. G. Bardakov, with proposed year 1999. The supplied record does not explain these differing attribution roles.
\item V. G. Bardakov, \emph{Algebra and Logic}, \textbf{39} (2000), 224--251, cited in the supplied statement for admissibility of every pair \((k,2)\) with \(k\geq2\). No article title or external link is supplied.
\end{itemize}
\end{document}
